\documentclass[11pt]{amsart}
\usepackage[T1]{fontenc}
\usepackage{lmodern,amsmath,amssymb,mathtools,microtype,booktabs,needspace}
\usepackage[margin=1.15in]{geometry}
\usepackage[hidelinks]{hyperref}
\newtheorem{theorem}{Theorem}[section]
\newtheorem{proposition}[theorem]{Proposition}
\newtheorem{lemma}[theorem]{Lemma}
\newtheorem{corollary}[theorem]{Corollary}
\theoremstyle{definition}\newtheorem{definition}[theorem]{Definition}
\theoremstyle{remark}\newtheorem{remark}[theorem]{Remark}
\newcommand{\R}{\mathbb R}
\newcommand{\Z}{\mathbb Z}
\newcommand{\E}{\mathbb E}
\newcommand{\Prob}{\mathbb P}
\newcommand{\dd}{\,\mathrm d}
\DeclareMathOperator{\covol}{covol}
\DeclareMathOperator{\area}{area}
\DeclareMathOperator{\rank}{rank}
\DeclareMathOperator{\dist}{dist}
\numberwithin{equation}{section}
\title[Reference-free certification]{Reference-free certification from intrinsic boundary laws}
\author{Qian Qi}
\date{September 29, 2026}
\subjclass[2020]{37D50, 53A04, 62L10, 62G05}
\keywords{Dispersing billiards, intrinsic boundary laws, completion certificates,
sequential discovery, unknown-cell sampling, period reconstruction}
\hypersetup{pdftitle={Reference-free certification from intrinsic boundary laws},pdfauthor={Qian Qi}}
\begin{document}
\begin{abstract}
We study completion and acquisition of intrinsic endpoint-law records for
periodic dispersing billiards. Two active endpoint density functions recover
a local reflecting pair and the absolute free-area normalization. The
nonnegative difference between an observed cycle covolume and the recovered
free and visible obstacle areas certifies both full coverage and period
saturation. A certificate alone does not acquire its records. We prove an
anytime-valid noisy stopping theorem for adaptively growing genuine-record
catalogues, with almost-sure finite stopping under divergent discovery
hazards and an explicit tail under uniform hazards. A randomized geometric
scanner realizes the hypotheses using stated body, clearance and periodic-gate
access, not unmarked trajectories alone. Uniform laboratory launches with
solid-position rejection approximate the unknown-cell Liouville law; their
boundary bias is charged through two growing endpoint histograms. The
results use bounded obstacle count, positive area and covolume gaps, analytic
strip, shape, symmetry and local physical margins. They are reference-free,
not prior-free. The local inverse, omission-safe certificate, reference-free
arithmetic and all earlier companion results are retained. No whole-table
minimax or exact analytic count-only conclusion is inferred.
\end{abstract}
\maketitle
\section{From a completion certificate to sequential acquisition}
\label{sec:setting}
A collection of local laws may reconstruct every obstacle it sees and
still omit a body orbit or a period. Absolute probability normalization
allows these omissions to be detected by one nonnegative defect. A second
question remains: when records arrive successively, with neither a known
catalogue nor a retained witness, can discovery and noisy certification
be combined into a sound stopping rule? We address this question while
separating the producer, the phase preparation and the geometric inverse.

\subsection{Physical records and the preparation contract}
Let
\[
 \mathcal O=\bigcup_{i=1}^r\bigcup_{\lambda\in\Lambda}(C_i+\lambda)
\]
be a periodic family of disjoint compact strictly convex planar bodies
with positive curvature and full-rank period lattice $\Lambda$. Its
free cell area is
\begin{equation}\label{eq:area-budget}
 A=\covol\Lambda-\sum_{i=1}^r\area(C_i)>0.
\end{equation}
A local record selects an isolated clear shortest bridge and its
three-impact return branch. The first and last impact positions are
measured by intrinsic arclength $s,t$ from the source normal contact.
One coherent reversal $(s,t)\mapsto(-s,-t)$ is allowed within a record;
origins and signs are not registered across records. At two known absolute
times the exact observation consists of two complete endpoint density
functions on a strictly active rectangle. Its finite counterpart consists
of two growing two-dimensional histograms. Failure and successful endpoints
outside that rectangle remain outcomes. These are not two scalar observations.

The ideal preparation has speed one and normalized cell phase measure
$\dd q\dd\vartheta/(2\pi A)$. It includes the free area affected by
unseen bodies. Conditioning on successful returns would destroy this
normalization. Section~\ref{sec:normalization} instead constructs it
approximately by uniformly drawing laboratory positions in large squares
and rejecting solid positions. No period vectors or cell boundary are
used in this preparation. Its finite-size error is retained in the proof.
The record evaluation rule must be periodic: it recognizes every translated
occurrence of the selected branch. A detector restricted to one physical
copy does not satisfy this contract.

The global class is analytic, asymmetric and shape separated: each $C_i$
has trivial Euclidean isometry group and different representatives are
noncongruent. Consequently $\Lambda$ is the full translation group of
$\mathcal O$. Quantitative results use known bounds on the number of
body types, diameters, lattice presentations, free area and cell covolume;
positive lower body-area, separation and curvature bounds; a common
analytic strip; shape and symmetry margins; and the local active, twist,
non-grazing, time-separation and onset-bracket bounds stated in
Section~\ref{sec:uniform}. Reference-free does not mean prior-free.

\subsection{Sequential discovery: the new conclusion}
A producer returns either no record or a replayable genuine local branch.
It may adapt to past observations, but validation uses fresh preparations.
It must supply the periodic evaluation gate and local validity, not a
claim that its catalogue is complete. Definition~\ref{def:producer}
formalizes this access. A finite saturated witness means a collection
covering every body type and generating the full period group; the witness
is unknown to the learner and appears only in the termination proof.

\begin{theorem}[Sequential discovery and certification]
\label{thm:discovery-main}
On the uniform physical class above, with periodic genuine-record gates,
there is a procedure receiving a growing, adaptively produced record list
with the following properties. For every prescribed accuracy $\xi>0$
and error probability $0<\delta<1/4$, all its whole-table certificates,
at all epochs and at data-dependent stopping times, are simultaneously
correct with probability at least $1-\delta$. An accepted reconstruction
has simultaneous support-function $C^2$ and period-basis error at most
$\xi$, modulo common isometry, permutation and representative choices.

If some finite saturated witness has divergent conditional discovery
hazards, the procedure stops almost surely. If its $w$ records each have
conditional proposal probability at least $p>0$ until discovered, then
for every sufficiently large deterministic epoch $k$,
\begin{equation}\label{eq:main-tail}
 \Prob\{T_\xi>k\}\le\frac{\delta}{8(k+1)^6}+w e^{-pk}.
\end{equation}
Its expected number of launch attempts is finite. No full record list,
reference table, witness, actual obstacle count or period vectors are
supplied to the stopping rule.

The preparation uses growing laboratory squares rather than an exact
unknown-cell sampler. With $N$ free preparations per histogram time,
$J$ records, bandwidth $h$ and square half-width $L$, the density error is
\begin{equation}\label{eq:main-bias}
 C\{h^\beta+\sqrt{r_N}\,h^{-3}+(r_N+L^{-1})h^{-4}\},\qquad
 r_N=\frac{\log(CJN/\alpha)}{N}.
\end{equation}
Here $0<\beta\le1$ and $C$ depends on all the stated priors. The
geometric scanner of Proposition~\ref{prop:geometric-producer} realizes
the discovery condition with explicit $w,p$ bounds. It uses body-list,
shortest-pair, clearance and periodic-gate access; this is not a theorem
that unmarked trajectories alone supply those primitives.
\end{theorem}

Sections~\ref{sec:normalization}--\ref{sec:sequential} prove this theorem.
Its two logically separate ingredients are probability control for every
validation epoch and a condition ensuring that sufficient genuine records
actually appear. The certificate alone does not acquire records. Without
the latter condition, it remains sound but need not stop. With geometric
access, the producer need not complete the whole catalogue before stopping.
The finite arithmetic stage is effective; compact analytic-body fitting
remains an existence construction. No whole-table minimax conclusion is
inferred from \eqref{eq:main-bias} or \eqref{eq:main-tail}.

\subsection{The deterministic inverse retained in the proof}
We state its local and global parts so the primary argument is self-contained.
Their full proofs occupy Sections~\ref{sec:local}--\ref{sec:uniform}.

\begin{theorem}[Intrinsic local inverse]\label{thm:local-main}
For a smooth isolated normal channel, two strictly active endpoint density
functions determine the two reflecting boundary germs in their relative
placement, the gap and the free area, up to a common Euclidean isometry.
The stationary two-flight action and area satisfy
\begin{equation}\label{eq:ratio-main}
 W=\frac{T_1f_2-T_2f_1}{f_2-f_1},\qquad
 A=\frac{(T_2-T_1)(-W_{st})}{2\pi(f_2-f_1)}.
\end{equation}
On the physical $C^3$ classes and all the positive margins stated in
Proposition~\ref{prop:local-stable}, this inverse is Lipschitz from the
two $C^2$ densities to the two intrinsic $C^2$ arcs. Analyticity and an
increasing jet order are unnecessary for this local assertion.
\end{theorem}

Recover the entire analytic pair for each record and identify visible
body types by congruence. Retain every record as a multigraph edge, including
loops and repetitions. For any connected observed component place body
representatives along a spanning tree. Its non-tree displacements $d_e$
are elements of the true period group. Set
\[
 I=\{\hbox{visible body types}\},\qquad
 S_I=\sum_{i\in I}\area(C_i),\qquad \Gamma=\sum_e\Z d_e.
\]
The construction does not use supplied integer copy labels.

\begin{theorem}[Exact whole-table completion certificate]
\label{thm:completion-main}
For any finite collection of genuine records of an analytic asymmetric
shape-separated periodic table, a connected rank-two component has defect
\begin{equation}\label{eq:defect-main}
 \mathcal D=\covol\Gamma-A-S_I
 =([\Lambda:\Gamma]-1)\covol\Lambda+\sum_{i\notin I}\area(C_i)\ge0.
\end{equation}
It vanishes precisely when the component sees every body orbit and
$\Gamma=\Lambda$. In this case it determines the entire periodic table
and its full unmarked period group. A positive defect or a rank-deficient
component is only a noncertificate. No coverage or supplied witness is
required for soundness, and arbitrary omissions or repetitions cannot
produce a false positive in the stated class.
\end{theorem}

\begin{theorem}[Reference-free finite-error certification]
\label{thm:uniform-main}
On the class specified in Section~\ref{sec:uniform}, there are $C,e_0>0$
and $0<\theta<1$ such that, for any finite list of genuine records with
simultaneous $C^2$ density error at most $e<e_0$, one can identify visible
incidence, recover the rational cycle arithmetic and estimate every
component's defect with error at most $Ce^\theta$. The gap
$g_0=\min(V_0,a_0)$ separates incomplete rank-two components from complete
saturated ones. On acceptance the entire table is recovered with error
at most $Ce^\theta$ in simultaneous body and lattice gauges. Constants
are uniform over repetitions and over the length of the record list.
No reference shapes, retained witness or integer edge labels are inputs.
The hypothesis concerns a finite genuine list, not a producer of that list.
\end{theorem}

The distinction between physical validity and empirical fit is essential.
A fit does not certify that an arbitrary gate is a genuine local branch.
The new sequential theorem preserves this distinction while allowing a
list to grow without a prescribed horizon. Earlier moment, registered,
relative-law and fixed-window testing results remain supplied in full as
companion material; Section~\ref{sec:comparison} explains their roles.

\section{The local action and two reflecting arcs}\label{sec:local}
\subsection{The physical return branch}
Place a normal contact pair at $(0,0)$ and $(g,0)$ for the proof only.
In transverse graph coordinates its boundaries are
$(-\psi_0(y),y)$ and $(g+\psi_1(y),y)$, where
$\psi_b(0)=\psi_b'(0)=0$ and $\psi_b''(0)=\kappa_b>0$.
The broken length has a nondegenerate minimum in its intermediate
coordinate: its second derivative there is $2(g^{-1}+\kappa_o)>0$.
The implicit-function theorem gives a unique stationary intermediate
impact for endpoints in a smaller fixed box. Strict channel clearance
identifies this branch with the physical specular trajectory.

Use intrinsic unit-speed boundary parametrizations $x(s)$ and $y(u)$,
with both tangents in the positive transverse direction at zero, and put
\begin{equation}\label{eq:action}
 D(s,u)=|x(s)-y(u)|,\qquad
 W(s,t)=D(s,u(s,t))+D(t,u(s,t)).
\end{equation}
The stationary equation is $D_u(s,u)+D_u(t,u)=0$. On a sufficiently
small rectangle $W$ is smooth, symmetric in its two arguments, and has
positive dispersing twist $-W_{st}>0$. At the normal orbit,
\[
 W(0,0)=2g,\qquad W_{st}(0,0)=-\frac1{2g(1+g\kappa_o)}.
\]
These statements hold uniformly on compact physical families with the
stated geometry and the required fixed derivative bounds.

\begin{lemma}[The residual-time density]\label{lem:density}
On a sufficiently short onset collar the endpoint density is
\begin{equation}\label{eq:density}
 f_T(s,t)=\frac{-W_{st}(s,t)}{2\pi A}(T-W(s,t))_+.
\end{equation}
On a rectangle where $T_1-W>0$, the two density functions at
$T_2>T_1$ recover $W$ and $A$ by \eqref{eq:ratio-main}.
\end{lemma}
\begin{proof}
Let $p$ be the momentum conjugate to source arclength and $r$ the time
before its first impact. Collision-section phase volume is
$\dd s\dd p\dd r$. The generating identity $p=-W_s$ changes this to
$(-W_{st})\dd s\dd t\dd r$. The allowed residual interval is
$0<r<T-W(s,t)$. Choose the collar shorter than a common positive lower
bound for adjacent free flights. Then this interval cannot include a
preceding collision, and the time after the third impact cannot include
a fourth. The strict local minimum and patch isolation keep the active
sublevel in this physical chart. Conversely, every selected near-onset
trajectory has these coordinates. Integrating $r$ and normalizing by
$2\pi A$ proves \eqref{eq:density}.

On a strictly active rectangle set $w=-W_{st}/(2\pi A)>0$.
Then $f_j=w(T_j-W)$ and $f_2-f_1=(T_2-T_1)w$.
Eliminating $w$ proves the action formula, and its recovered mixed
second derivative proves the area formula. Equality of measures
uniquely determines their continuous density representatives.
\end{proof}

\begin{proposition}[Exact local information equivalence]\label{prop:lens}
On the physical class and a fixed strictly active rectangle, the two
endpoint density functions at known absolute times are equivalent to
$(W,A)$. In particular they determine the local canonical scattering
relation
\begin{equation}\label{eq:lens}
       (s,-W_s(s,t))\longmapsto(t,W_t(s,t)).
\end{equation}
Conversely, that relation, one absolute action value, and $A$ determine
the densities on the rectangle. No extension to a global exterior
travelling-time relation is asserted.
\end{proposition}
\begin{proof}
The forward and inverse formulas are Lemma~\ref{lem:density}.
The positive twist makes $t\mapsto-W_s(s,t)$ locally invertible.
Thus \eqref{eq:lens} is a canonical graph; its generating one-form is
$-p_-\dd s+p_+\dd t=\dd W$. On the connected rectangle it determines
$W$ up to a constant, fixed by one value. Substitution in
\eqref{eq:density} gives the converse. Here the given relation is assumed
physical, so its generating one-form is exact. Arbitrary positive
functions satisfying the quotient formula are not being declared
physical endpoint laws.
\end{proof}

This makes explicit both the richness and the origin of the data.
First derivatives of travel time recover endpoint directions in
scattering geometry; compare Noakes--Stoyanov~\cite[Lemma 4]{NS} and
Gurfinkel--Noakes--Stoyanov~\cite[Lemma 7]{GNS}. The affine elimination
is not a new abstract lens-rigidity principle. What remains to prove
here is an explicit inverse for two unknown reflecting arcs in their
intrinsic coordinates and its periodic descent with volume calibration.

\subsection{The curvature identities}
On the diagonal write $u(s)=u(s,s)$ and define
\begin{equation}\label{eq:diagonal}
 \ell(s)=\tfrac12W(s,s),\qquad a(s)=W_s(s,s),\qquad
 v(s)=\sqrt{1-a(s)^2}.
\end{equation}
Shrink the interval so that $v>0$. Let $T=x'$ and $N$ be the source
unit tangent and inward normal, with $T'=kN$ and $N'=-kT$.
The vector $e=(x-y)/\ell$ then satisfies $e=aT+vN$.

\begin{lemma}[Diagonal Schur complement]\label{lem:curvature}
The source curvature and the opposite curvature at the nearest point are
\begin{align}
 k(s)&=\frac{W_{ss}(s,s)-W_{st}(s,s)-v(s)^2/\ell(s)}{v(s)},
                                                    \label{eq:curvature-source}\\
 k_o(u(s))&=\frac1{\ell(s)}
       \left(\frac{v(s)^2}{-2\ell(s)W_{st}(s,s)}-1\right),
                                                    \label{eq:curvature-other}\\
 u'(s)&=\frac{v(s)}{1+\ell(s)k_o(u(s))}>0.          \label{eq:foot-speed}
\end{align}
\end{lemma}
\begin{proof}
At $s=t$ the stationary equation is $D_u=0$, selecting the nearest
point on the opposite strictly convex patch. Orient its tangent $T_o$
continuously so that $T\cdot T_o=v>0$. Direct differentiation gives
\[
 D_{ss}=v^2/\ell+kv,\qquad D_{su}=-v/\ell,\qquad
 D_{uu}=1/\ell+k_o.
\]
For the last identity, $y''=-k_oe$ at the nearest point; this is the
external, dispersing sign. The stationary elimination of the intermediate
variable gives
\begin{equation}\label{eq:schur}
 W_{ss}=D_{ss}-\frac{D_{su}^2}{2D_{uu}},\qquad
 W_{st}=-\frac{D_{su}^2}{2D_{uu}}
        =-\frac{v^2}{2\ell(1+\ell k_o)}.
\end{equation}
Subtracting the first two expressions gives
\eqref{eq:curvature-source}; solving the last gives
\eqref{eq:curvature-other}. Differentiation of $D_u(s,u(s))=0$
yields \eqref{eq:foot-speed}.
\end{proof}

\begin{proof}[Proof of Theorem~\ref{thm:local-main}, exact assertion]
Recover $W$ and $A$ from Lemma~\ref{lem:density}. Choose
$x(0)=0$, $T(0)=(0,1)$ and $N(0)=(-1,0)$ and solve the Frenet
system with the recovered $k$. This determines the source arc. The
opposite arc in the same plane is
\begin{equation}\label{eq:foot}
 y(u(s))=x(s)-\ell(s)\{a(s)T(s)+v(s)N(s)\}.
\end{equation}
By \eqref{eq:foot-speed} it is a regular arc containing
$y(0)=(g,0)$, where $g=W(0,0)/2$. The recovered objects are functions
on intervals, not merely their Taylor series. Simultaneous arclength
reversal reflects the entire pair, not its members separately. The
remaining choices are one translation and orthogonal transformation.
\end{proof}

\begin{proposition}[Conditional local Lipschitz inverse]\label{prop:local-stable}
Fix nested endpoint rectangles $Q_0\Subset Q$, a diagonal subinterval,
and positive constants $M,g_*,A_*,v_*,w_*,\tau_*,m_*$. Consider smooth
physical pairs for which, in their normal contact frames, both intrinsic
boundary parametrizations have $C^3$ norm at most $M$ on fixed slightly
larger arcs, $g_*,A_*\le g,A\le M$, $|T_j|\le M$,
$T_2-T_1\ge\tau_*$, $T_1-W\ge m_*$ on $Q$,
$-W_{st}\ge w_*$ on $Q$, $v\ge v_*$ on the diagonal, and
$\|f_j\|_{C^2(Q)}\le M$. The patches, isolation and adjacent free-flight
margins are uniform, as is positive curvature. For two such pairs at
the same times put
\[
        e=\max_{j=1,2}\|f_j-\widetilde f_j\|_{C^2(Q)}.
\]
After a permitted common reflection, their gaps, areas, and intrinsic
$C^2$ boundary arcs on smaller fixed intervals differ by at most $Ce$.
The constant depends on all the stated bounds and margins, not on a
contact-jet order. This is an estimate between physical density pairs,
not an inverse defined on every pair of $C^2$ functions.
\end{proposition}
\begin{proof}
The density difference is bounded below by
$\tau_*w_* /(2\pi M)$. Products and reciprocals are Lipschitz on the
corresponding bounded $C^2$ sets. Formula~\eqref{eq:ratio-main} gives
$C^2$ action error $Ce$; its area expression gives scalar error $Ce$.
The diagonal formulas give $Ce$ source-curvature error and opposite
curvature error as functions of $s$, since all their denominators have
positive margins. The Frenet integral equations and Gronwall's inequality
give $Ce$ error for the source arc through its second derivative.

For the opposite arc, differentiating \eqref{eq:foot} twice would
unnecessarily require a third derivative of $W$. Instead integrate
\eqref{eq:foot-speed}. Its speed has positive upper and lower bounds,
and its error is $Ce$. Consequently $u(s)$ and its inverse have error
$Ce$ on smaller intervals. The physical $C^3$ prior bounds the derivative
of opposite curvature in intrinsic arclength. Thus the two opposite
curvature functions differ by $Ce$ at the same $u$. Their Frenet systems,
with the recovered gap and common initial tangent, give the required
$C^2$ estimate. This uses exactly the physical prior stated in the
proposition and no derivative of an arbitrary noisy curvature estimate.
\end{proof}

\section{Observed cycles and the missing-area identity}\label{sec:completion}
\subsection{Complete images and visible incidence}
\begin{lemma}[Analytic images and unique matching]\label{lem:images}
A nonempty open arc of a compact connected embedded real-analytic
boundary determines its complete image among such boundaries. If $C,D$
are congruent asymmetric compact bodies, exactly one Euclidean isometry
maps $C$ onto $D$.
\end{lemma}
\begin{proof}
For the first assertion consider the points where two candidate curves
have the same germ. This set is open. At a limit point, local analytic
graph representations and the one-variable identity theorem extend
equality of the approaching common arcs, so the set is also closed.
Connectedness gives equality of the complete curves. The convex body
bounded by the curve is then fixed. For the second assertion, the
quotient of any two matching isometries is a symmetry of $C$ and is
therefore the identity.
\end{proof}

The local inverse and Lemma~\ref{lem:images} reconstruct each pair in
its relative placement. Shape separation makes congruence of its
individual members identify precisely their true obstacle types. This
gives the finite visible multigraph. Repeated experiments remain parallel
edges; an edge joining two copies of the same type becomes a loop.
We do not infer identities by comparing the experimental time windows.

\begin{lemma}[Assembly of an arbitrary connected observed component]
\label{lem:assembly}
For any connected component of this multigraph choose a root and a
spanning tree. The records determine one representative of every visible
body in a common plane. Each non-tree edge determines a displacement
$d_e\in\Lambda$. The group $\Gamma=\sum_e\Z d_e$ is independent of
the spanning tree and of all representative and reversal choices, up to
the orthogonal part of one ambient Euclidean isometry.
\end{lemma}
\begin{proof}
Place the root body arbitrarily. For an outward tree edge, match the
source member of its reconstructed pair to the already placed body.
Lemma~\ref{lem:images} gives a unique map; applying that same map to
the partner places the child. In the true periodic table these choices
amount to choosing one lifted representative of each visible type, with
tree labels zero. The labels are used only to justify the construction,
not supplied to it.

For a non-tree edge perform the same source matching. Its target image
$\widehat C_{j,e}$ and the placed target representative $C_j$ are
translates, because their true types agree. This translation is unique;
using the Steiner point $s(C)$ it is
\[
 d_e=s(\widehat C_{j,e})-s(C_j),\qquad
                         \widehat C_{j,e}=C_j+d_e.
\]
It lies in the true period lattice. Reversing a local arclength choice
reflects the whole pair and composes the matching map with that reflection,
leaving its placement unchanged. Reversing an edge changes the sign of
its cycle translation. Changing representative copies adds a vertex
coboundary to the edge translations; these differences cancel on closed
walks. Fundamental cycles generate all closed walks integrally. Thus
the group does not depend on the tree. Parallel repetitions either
repeat a generator or introduce a zero cycle, and cause no problem.
\end{proof}

\subsection{The certificate and its consequences}
\begin{proof}[Proof of Theorem~\ref{thm:completion-main}]
The component identifies exactly a subset $I$ of the true types, because
of shape separation, even though $r$ and the complement of $I$ are
unknown. Lemma~\ref{lem:assembly} gives $\Gamma\subset\Lambda$.
If its rank is two, the index $m=[\Lambda:\Gamma]$ is a positive
integer and $\covol\Gamma=m\covol\Lambda$. Subtracting the
normalization identity \eqref{eq:area-budget} gives
\begin{equation}\label{eq:defect}
 \covol\Gamma-A-S_I=(m-1)\covol\Lambda+S_{I^c},\qquad
 S_{I^c}=\sum_{i\notin I}\area(C_i).
\end{equation}
Both terms on the right are nonnegative. Every obstacle has positive
area. Hence equality to zero is equivalent to $m=1$ and $I^c$ empty.
In that case all representatives have been placed and the full period
group is $\Gamma$, which determines the table. Any other physical
table in the stated class giving these records must have the same
visible images, cycle group and $A$. Applying \eqref{eq:defect} to it
also excludes hidden types and a larger period group. Thus the inverse
is unique even among candidates not initially assumed to be covered.
\end{proof}

\begin{corollary}[A quantitative separation of incomplete data]
\label{cor:defect-gap}
Suppose $\covol\Lambda\ge V_0>0$ and each obstacle orbit has area at
least $a_0>0$. Every connected rank-two observed component has either
$\mathcal D=0$ or
\begin{equation}\label{eq:defect-gap}
                    \mathcal D\ge g_0:=\min(V_0,a_0).
\end{equation}
\end{corollary}
\begin{proof}
If $m>1$, the first term of \eqref{eq:defect} is at least $V_0$.
If some type is unseen, the second term is at least $a_0$.
\end{proof}

\begin{corollary}[Finite positive stopping certificate]\label{cor:stopping}
Any acquisition source returning genuine records may be stopped once
some connected rank-two component has zero defect. No prior claim
that the source retained a specified witness is needed for soundness.
If the accumulated records eventually include a connected covering set
whose cycles generate $\Lambda$, this stopping criterion is eventually
satisfied. Arbitrary deletions can prevent acceptance but cannot turn
an incomplete physical component into an accepted one.
\end{corollary}
\begin{proof}
Soundness is Theorem~\ref{thm:completion-main}. Adding records preserves
a previously covered saturated component and can only join visible
components or add generators. A covering saturated set gives zero
defect when all of its finitely many records are present. Apply the same
theorem to any retained subcollection for the last assertion.
\end{proof}

The converse does not say that every useful list is accepted: a disconnected
list can contain much geometric information but does not register all
its components. The criterion tests components separately. Nor does
nonacceptance identify which acquisition mechanism failed. It simply
withholds a whole-table assertion.

\subsection{Arithmetic without assuming coverage}
For exact displacements pick any independent pair $D=[d_1\ d_2]$.
There is an integer $n_D=|\det D|/\covol\Lambda$ such that every
coordinate vector $z_e=D^{-1}d_e$ belongs to $n_D^{-1}\Z^2$.
This follows because the finite group $\Lambda/D\Z^2$ has order
$n_D$, hence is annihilated by $n_D$. Let $q$ be a common denominator
of the finitely many $z_e$; the least one divides $n_D$. Form the
integer matrix with columns $qz_e$, including $qe_1,qe_2$. If $H$
is a column basis for their subgroup of $\Z^2$, then
\begin{equation}\label{eq:denominator-group}
 \Gamma=\frac1qDH\Z^2,\qquad
 \covol\Gamma=\frac{|\det D|\,|\det H|}{q^2}.
\end{equation}
An elementary Hermite reduction computes $H$. Equivalently its determinant
is the gcd of the two-column minors: factor the integer column matrix
as $HE$, note that the columns of $E$ generate $\Z^2$, and apply
Cauchy--Binet to an integer right inverse of $E$. This gives gcd one
for its minors and therefore gcd $|\det H|$ for the original matrix.
The general algorithms are classical \cite{KB}; the rank-two argument
here is included to fix the column convention.

Crucially, $n_D$ is not computed as $|\det D|/(A+S_I)$: that substitution
would already assume the conclusion that $I$ contains every type.
Section~\ref{sec:uniform} recovers the rational coordinates using only
an upper denominator bound and proves that they can be identified under
noise before the completion test is run.

\subsection{Physical examples with unobserved obstacles}
\begin{proposition}[Both defects occur in analytic tables]\label{prop:examples}
There are analytic asymmetric shape-separated periodic tables and clear
local records realizing each of the following cases: full cycle generation
with an unseen obstacle; complete shape coverage with period index two;
and zero defect without a primitive pair of recorded cycles.
\end{proposition}
\begin{proof}
Work on $\Lambda=\Z^2$ with a sufficiently small asymmetric analytic
body at every lattice point. For example a small multiple of the support
function
\[
 p(\varphi)=1+\epsilon\cos(2\varphi)
                   +\epsilon\cos(3\varphi+\pi/4),\quad
                         0<|\epsilon|<1/11,
\]
is strictly convex since $p+p''>0$. A rotation preserving its nonzero
second and third modes is trivial. A reflection with axis angle $a$
would require both $2a\in\pi\Z$ and $3a+\pi/4\in\pi\Z$, which
is impossible. It is therefore asymmetric.

The shortest bridges associated with the primitive vectors
$(1,0),(1,2),(-1,3)$ are clear for small bodies: the corresponding
center segments contain no other lattice points and have positive
clearance from all other centers. Their pair determinants have absolute
values $2,3,5$, whose gcd is one, so the three cycles generate $\Z^2$
although no pair does. This gives zero defect. Keeping only $(1,0)$
and $(1,2)$ gives index two and defect one, while still observing the
only shape.

Choose a point on the torus away from the finite union of all these
segments and their short experimental tubes, and put a second, differently
scaled asymmetric body there and at all its lattice translates. It can
be made small enough to leave every selected return branch unchanged.
The representatives are noncongruent. Its area $a$ reduces the free area
by exactly $a$. If only the first type is observed using all three
cycles, the cycle group is $\Z^2$ and the defect is $a$, not zero.
With only the first two cycles the defect is $1+a$. These are actual
periodic tables, not only formal area budgets. Strict clearances and
shape inequalities persist under small suitable perturbations.
\end{proof}

\section{A noisy certificate without a reference table}\label{sec:uniform}
The two difficulties in applying \eqref{eq:defect} to noisy data are
identifying body types without a known witness and recovering a discrete
subgroup before its covolume has been calibrated. The first requires
only shape separation, not isolation of every edge key. For the second
we use bounded-denominator separation and postpone the area test until
the subgroup has already been recovered.

\subsection{Uniform physical class}
Fix an upper bound $r_0$ for the number of obstacle types, a diameter
bound $D_0$, and a bound $R$ for every recorded gap. The actual number
and identities of types are unknown. Require $A_0\le A\le A_1$, a
positive lower bound $V_0$ for the true cell covolume, and a positive
lower bound $a_0$ for every obstacle area. One may take $V_0=A_0$;
we keep the symbol separate to display its arithmetic role. There are
uniform positive separation, curvature, channel-isolation and adjacent
free-flight bounds. Lattice bases and their inverses admit uniform bounds
in some reduced presentation. Each observed record satisfies the physical
$C^3$, active, twist, diagonal non-grazing, time-separation and rectangle
conditions of Proposition~\ref{prop:local-stable}, and its densities
have a common $C^{2,\beta}$ bound on a slightly larger rectangle.

After Steiner centering, each support function $p$ is holomorphic and
bounded by $M$ on $|\operatorname{Im}\varphi|<\rho$, for fixed
$M,\rho>0$. Local contact graphs have common bounded holomorphic
extensions on a fixed complex disk. These graph bounds also follow
on smaller disks from the support bounds and positive curvature.
For positive $\eta,\zeta$ require
\begin{align}
 \|p(\cdot+\phi)-p\|_2&\ge
            \eta\dist(\phi,2\pi\Z),\label{eq:rotation}\\
 \|p(-\cdot+\phi)-p\|_2&\ge\eta,\label{eq:reflection}\\
 \inf_{O\in O(2)}\|p_i-p_{OC_j}\|_2&\ge\zeta\quad(i\ne j).
                                                   \label{eq:shape-separation}
\end{align}
Take a closed class with these margins, compact after restricting to a
narrower strip and bounded root placement. Bounds and margins, but no
reference shape, period basis, actual obstacle count or reference witness,
are inputs. The class need not lie in one local deformation chart.
The margins allow infinite-dimensional analytic shape variation.

All comparisons below allow one common Euclidean motion, a permutation
of obstacle types, and a change of representative copies and period basis.
A table error at most $t$ means that such representatives and bases can
be chosen so that all support-function $C^2$ differences and the period
matrix difference are at most $t$. We prove estimates in these explicit
simultaneous gauges; no continuous choice of a canonical lattice basis
across its reduction walls is asserted.

\subsection{From densities to complete pair images}
\begin{lemma}[Uniform pair recovery]\label{lem:pair-error}
If fitted physical density pairs have $C^2$ error at most $e<1$ on
the common rectangles, each recovered entire analytic pair in its own
contact frame has support-function $C^2$ error at most
$\Delta=C e^\theta$, where $C>0$ and $0<\theta<1$ depend only
on the stated class. The recovered free area has error at most $Ce$.
\end{lemma}
\begin{proof}
Proposition~\ref{prop:local-stable} gives $Ce$ errors for the two
intrinsic arcs and $A$. Positive curvature converts a smaller arc to
a normal-angle interval of a fixed positive length. The inverse angle
map and its support values depend locally Lipschitzly on the $C^2$ arc,
using the uniform physical $C^3$ bound. Hence the support functions
are $Ce$-close on a fixed real interval in the contact frame.

Here is the continuation estimate. Slit a narrower periodic strip
cylinder along a shorter observed interval. The harmonic function equal
to one on both slit sides and zero on the two outside boundary circles
is positive in the slit cylinder. On a still narrower closed cylinder,
including the slit by continuity, it has a positive minimum $\theta$.
The endpoints of the slit are regular Dirichlet boundary points.
Apply the subharmonic maximum principle to the logarithm of the modulus
of the support difference, using $Ce$ on the slit and its uniform
complex bound elsewhere. This gives $C'e^\theta$ on the narrower
cylinder. Zeros follow by regularization; Cauchy estimates on fixed
disks give the same bound through two derivatives on the real circle.
Shifting from centered to contact frames changes the complex support
bound by a controlled diameter term. Constants are uniform under rotations
and under translating the observed interval. This is conditional analytic
continuation of functions on arcs, not an infinite jet estimate; compare
\cite{Trefethen}.
\end{proof}

Fitting can be defined independently for each record over a fixed compact
physical pair class. Given a density estimate with certified error $e$,
choose a point of a fixed countable dense family whose $C^2$ discrepancy
is within $e$ of the infimum. It then differs from the true density by
at most $3e$. A first-eligible-point rule is measurable. This is an
existence construction, not an efficient analytic continuation algorithm.
It does not declare every noisy pair of functions to be a physical law.
Uniform physical branch validity remains a model assumption.

\begin{lemma}[Incidence and placements without edge-key neighborhoods]
\label{lem:reference-free-assembly}
For sufficiently small $\Delta$, the recovered body shapes identify
all visible types without a reference table. A spanning tree chosen in
the resulting observed multigraph gives representative and cycle-vector
errors at most $C\Delta$, uniformly over repetitions and surplus edges.
No classification of distinct channel orbits is required.
\end{lemma}
\begin{proof}
Use the centered support distance modulo $O(2)$ in $L^2$. Two estimates
of the same body type are at distance at most $C\Delta$; estimates
of different true types are at distance at least $\zeta-C\Delta$.
For $C\Delta<\zeta/4$, joining observations whose distance is less
than $\zeta/2$ gives exactly the true type clusters. There is no
ambiguous chain of clusters, since the two cases have disjoint distance
ranges. Each record is retained as one edge with its two recovered
endpoint types. Multiple edges and loops are allowed, so near-coincidence
of two different edge geometries never has to be resolved.

For a tree placement, choose a least-discrepancy orthogonal match of
centered source shapes and then match their Steiner points. Compactness
of $O(2)$ ensures existence. A true matching has discrepancy $O(\Delta)$,
so the minimizing match does too. Relative to the true matching its
orthogonal action moves the true support by $O(\Delta)$ in $L^2$.
Equation~\eqref{eq:reflection} excludes the wrong component, and
\eqref{eq:rotation} bounds the remaining rotation angle by
$O(\Delta/\eta)$. The uniform $C^3$ support bound controls the
resulting $C^2$ placement error. Steiner-point translations are bounded
linear functionals of support values. Induction along a tree path of
at most $r_0-1$ edges proves the representative bound.

For every non-tree edge, match its source in the same way and subtract
the Steiner point of the placed target. The corresponding exact physical
operation gives $d_e\in\Lambda$, and the displacement error is
$O(\Delta)$. Constants depend on $r_0$ and the margins, not the
number of repeated or unused records. This comparison uses the truth
only in the proof; the choice of root, tree and matches is entirely
from the fitted observations.
\end{proof}

\subsection{Rational reconstruction before calibration}
Set
\begin{equation}\label{eq:denominator-bound}
 B_0=(2r_0-1)(R+2D_0),\qquad
                  Q=\max\{1,\lceil B_0^2/V_0\rceil\}.
\end{equation}
Every exact non-tree displacement has norm at most $B_0$. Indeed
representative Steiner points in a rooted tree lie within
$(r_0-1)(R+2D_0)$ of the root; an edge pair has Steiner-point distance
at most $R+2D_0$. Subtracting the placed target gives the stated bound.

\begin{lemma}[Reference-free discrete recovery]\label{lem:rational}
Suppose estimated displacements have uniform error at most $t$ and
$t$ is sufficiently small in terms of $B_0,V_0$. One can determine
whether their true rank is two. In the rank-two case choose a pair
$D$ by maximizing the estimated absolute determinant. Each coordinate
of $D^{-1}d_e$ is a rational with denominator at most $Q$. For all
sufficiently small $t$, these coordinates, their least common denominator
$q$, and the integer column group in \eqref{eq:denominator-group} are
recovered exactly by a finite rational search, using no area calibration.
Moreover $q\le Q$, and the estimated cycle covolume has error at most
$Ct$.
\end{lemma}
\begin{proof}
The determinant error of each pair is at most $(2B_0+t)t$. In a true
rank-two component, every nonzero pair determinant is an integer multiple
of $\covol\Lambda$ and hence at least $V_0$. In a rank-one component
all pair determinants vanish. Once the error is less than $V_0/4$,
a maximal estimated determinant exceeding $V_0/2$ selects an independent
true pair, and no dependent component passes this test. The selected
true determinant lies between $V_0$ and $B_0^2$.

Write $n_D=|\det D|/\covol\Lambda\in\{1,\ldots,Q\}$. The
finite-group argument preceding \eqref{eq:denominator-group} puts each
coordinate in $n_D^{-1}\Z$. Its absolute value is bounded, for example,
by $2B_0^2/V_0$. Matrix inversion on this determinant-separated set
gives coordinate error at most $C_0t$. Distinct reduced rationals with
denominators at most $Q$ differ by at least $Q^{-2}$: their difference
has a nonzero integer numerator over a denominator at most $Q^2$.
If $C_0t<1/(4Q^2)$, the true value is the unique rational in the
allowed bounded range within $1/(4Q^2)$ of the estimate. A finite
search over denominators and bounded numerators recovers it exactly.
All actual denominators divide the same $n_D$, so their least common
multiple $q$ divides $n_D$ and is at most $Q$.

Hermite reduction now gives the same integer column group and a basis
$H$ for the truth and its estimate. Since this group contains
$q\Z^2$, it has $1\le|\det H|\le q^2$. Formula
\eqref{eq:denominator-group} with the estimated $D$ gives the covolume
estimate; its error is at most the pair-determinant error. All computations
of integer structure precede any use of $A+S_I$.
\end{proof}

This is not integer spanning of noisy real vectors, which could give
a dense subgroup. The finite denominator bound comes from periodicity
and a covolume lower bound; rational separation fixes the exact arithmetic
first. The search may be large, but is finite once the stated priors
are fixed. It does not make the preceding physical fitting efficient.

\subsection{The decision and reconstruction}
After type clustering, consider each connected component. Reject rank
less than two. In rank two compute $q,H$ as above, and put
\begin{equation}\label{eq:estimated-defect}
 \widehat{\mathcal D}
 =\frac{|\det\widehat D|\,|\det H|}{q^2}
           -\widehat A-\sum_{i\in I}\area(\widehat C_i).
\end{equation}
The estimate $\widehat A$ may be taken from any one fitted record;
all such estimates are uniformly close to the common true value. Use
one fitted placed body per type, rather than summing repetitions.

\begin{proof}[Proof of Theorem~\ref{thm:uniform-main}]
Lemmas~\ref{lem:pair-error}--\ref{lem:rational} give the correct visible
clusters, graph and rational cycle structure when $e<e_0$ for a
class-dependent threshold. For a $C^2$ support function,
\[
 \area(C)=\tfrac12\int_0^{2\pi}(p^2-(p')^2)\dd\varphi.
\]
On the bounded class this is Lipschitz in the $C^1$ norm. There are at
most $r_0$ visible types, so \eqref{eq:estimated-defect} differs from
\eqref{eq:defect} by at most $C\Delta$. Reduce $e_0$ until this is
less than $g_0/4$, with $g_0=\min(V_0,a_0)$. Accept precisely when
\begin{equation}\label{eq:acceptance}
                     |\widehat{\mathcal D}|<g_0/2.
\end{equation}
Corollary~\ref{cor:defect-gap} shows that every complete saturated
component passes and every incomplete component fails. Neither witness
retention nor coverage is assumed in proving soundness of the decision.

On acceptance all types are visible and $\Gamma=\Lambda$. The
reconstructed period basis is $\widehat B=\widehat D H/q$, while
$B=DH/q$ is a basis of the true period lattice. The integer group
contains $q\Z^2$; there are only finitely many such column Hermite
forms for $q\le Q$, so $\|H/q\|$ has a bound depending only on $Q$.
Thus $\|\widehat B-B\|\le C_Q\Delta$. The placed body supports
are simultaneously $C\Delta$-close by Lemma~\ref{lem:reference-free-assembly}.
This gives the stated whole-table bound in common representative and
period gauges. Small errors and the strict curvature and global separation
margins keep this finite representation a physical periodic table:
bounded bases and their inverses reduce potentially close translated
pairs to a uniformly finite set; the remaining pairs are farther apart.
Alternatively one can project to the compact admissible physical class
with an additional factor in the error, since the true table is within
that error of the construction. A countable dense near-minimizer gives
a measurable selection. All decisions use the observations and global
bounds, never a reference table or local witness-key neighborhoods.
\end{proof}

The decision is uniform over any data-dependent deletion of records on
a simultaneous accuracy event for the original finite list. Surplus
channels need no common key separation and may change as long as the
ones actually used remain genuine branches within the local physical
class. No theorem is asserted for fabricated records or an unrestricted
mixture of unseparated itineraries. A failed test is an informative
noncertificate, not evidence that a table does not exist.

\section{Phase preparation without a known fundamental cell}
\label{sec:normalization}
The completion defect uses the free area of the \emph{whole} quotient.
A probability conditional on a successful return, or on launching near
one selected channel, cannot replace it. We now specify a preparation
which needs neither the period vectors nor the unknown obstacle list.
At finite scale it is approximate, and its error will remain in every
subsequent statistical bound.

\subsection{A periodic gate and a laboratory launch}
A record gate is a measurable function of a launched trajectory, returning
its two intrinsic endpoint coordinates on the selected branch, or a
failure symbol. It is \emph{periodic} if translating a launch and its
trajectory by any $\lambda\in\Lambda$ leaves the output unchanged.
Thus it counts all translates of its selected local branch by the same
local rule. It supplies no integer copy label. The local gate, its
paired-window identity and its coherent arclength reversal are retained
sensor capabilities. A gate addressed to one physical copy in the plane
is not periodic and is not covered by this construction.

The apparatus draws $q$ uniformly in the known laboratory square
$\Omega_L=[-L,L]^2$ and tests only whether $q$ is free. Solid launches
are rejected before a trajectory is launched. On a free launch it draws
a uniform direction, follows the billiard, and retains every gate outcome,
including failure. The apparatus need not report a geometric map of its
rejected positions. This is rejection in the \emph{launch location}, not
conditioning on success of a collision itinerary. Both rejected launch
attempts and all accepted free preparations are counted below.

Assume known bounds $A\ge A_0>0$, $V:=\covol\Lambda\le V_1$, and
that some fundamental parallelogram has diameter at most $b$. Its vectors,
position and free subset are not known. Such bounds follow from the
bounded lattice presentation in Section~\ref{sec:uniform}; for instance
$V_1=A_1+r_0\pi D_0^2$ is a harmless upper bound. Write $\nu_L$
for the projection modulo $\Lambda$ of an accepted free launch and its
direction, and $\mu$ for normalized quotient Liouville measure. Total
variation is $\|\nu-\mu\|_{\rm TV}=\sup_E|\nu(E)-\mu(E)|$.

\begin{proposition}[Uniform cell-free approximation]\label{prop:box}
For $L\ge2b$,
\begin{equation}\label{eq:box-tv}
 \|\nu_L-\mu\|_{\rm TV}\le
 \tau_L:=\min\left\{1,\frac{8V_1b}{A_0L}\right\}.
\end{equation}
The same bound holds after every periodic measurable record gate, for
every observation time and every cell of an endpoint histogram. The
free-launch probability is at least $p_0=A_0/(4V_1)$. Producing $N$
accepted free preparations takes at most $N/p_0$ launch attempts in
expectation, and at most
\begin{equation}\label{eq:launch-cost}
 \left\lceil\frac{8\{N+\log(1/\alpha)\}}{p_0}\right\rceil
\end{equation}
attempts with probability at least $1-\alpha$.
\end{proposition}
\begin{proof}
Tile the plane by a half-open fundamental parallelogram of diameter at
most $b$. Let $U$ be the union of tiles wholly contained in $\Omega_L$.
Every point of $[-L+b,L-b]^2$ lies in such a tile. Consequently
\[
 |\Omega_L\setminus U|\le8Lb,\qquad |U|\ge4(L-b)^2.
\]
If $U$ contains $n$ tiles, its free part has area $nA=(A/V)|U|$.
Conditional on a free launch in $U$, projection is \emph{exactly}
uniform on the free quotient; no knowledge of this tiling is used by
the apparatus. The projected measure from the remaining free part is
arbitrary, but its mixture weight is at most
\[
 \frac{8Lb}{(A_0/V_1)4(L-b)^2}\le\frac{8V_1b}{A_0L}.
\]
This proves \eqref{eq:box-tv}. Independent uniform direction and a
measurable gate cannot increase total variation. The free area in $U$
also gives launch acceptance probability at least
$(A_0/V_1)(1-b/L)^2\ge p_0$.

The number of attempts to get $N$ acceptances is a sum of independent
geometric waiting times, giving the expectation. For $m$ attempts its
upper tail is the probability that a binomial variable of mean at least
$mp_0$ is below $N$. If $m$ is the ceiling in \eqref{eq:launch-cost},
then $N\le mp_0/8$. The multiplicative Chernoff bound makes this
probability at most $\exp(-mp_0(7/8)^2/2)\le\alpha$.
\end{proof}

This argument is a direct periodic boundary-layer estimate, not an
assumption of mixing or equidistribution along one billiard orbit. It
remains valid in the presence of infinite horizons. Unknown hidden
bodies alter the free-launch distribution itself, so their area is not
removed by any later aperture normalization.

\subsection{Propagation through a finite histogram}
Fix nested active endpoint rectangles as in Proposition~\ref{prop:local-stable}.
Use a square mesh of side $h$, including a surrounding collar of cells.
The grid is centered so simultaneous arclength reversal permutes cells.
Let $f$ be a true cell-normalized density with a common
$C^{2,\beta}$ bound, $0<\beta\le1$. Cell probabilities satisfy
$p_B\le C h^2$. The large-square experiment instead has probabilities
$p_{B,L}$ with $|p_{B,L}-p_B|\le\tau_L$.

\begin{lemma}[Histogram estimate with preparation bias]
\label{lem:biased-histogram}
For $J$ fixed genuine records, $N$ independent accepted free preparations
at each of the two histogram times, and failure probability $\alpha$,
put $r_N=\log(CJN/\alpha)/N$. For $r_N<1$ and bandwidths with
$h^{-2}\le N$, a local quadratic reconstruction satisfies, simultaneously
for all records and both densities,
\begin{equation}\label{eq:biased-rate}
 e_N(h,L):=\|\widehat f-f\|_{C^2}
 \le C\{h^\beta+\sqrt{r_N}\,h^{-3}
                       +(r_N+\tau_L)h^{-4}\}
\end{equation}
with probability at least $1-\alpha$. In particular, choosing
\begin{equation}\label{eq:box-bandwidth}
 h=r_N^{1/(2\beta+6)},\qquad
 L\ge C r_N^{-(\beta+4)/(2\beta+6)}
\end{equation}
gives $e_N\le C r_N^{\beta/(2\beta+6)}$. The constants depend on
the physical and preparation bounds, not on a known cell or a known list
of obstacle shapes.
\end{lemma}
\begin{proof}
For one cell, Bernstein's inequality and
$p_{B,L}\le Ch^2+\tau_L$ bound its sampling error, after a union bound,
by $C\{h\sqrt{r_N}+\sqrt{\tau_Lr_N}+r_N\}$. Adding its deterministic
bias and using $\sqrt{\tau_Lr_N}\le(\tau_L+r_N)/2$ gives
\begin{equation}\label{eq:cell-error}
 |\widehat p_B-p_B|\le C\{h\sqrt{r_N}+r_N+\tau_L\}.
\end{equation}
There are at most $CJh^{-2}\le CJN$ cell events, so the stated logarithm
suffices. Different cell indicators from one launch need not be independent;
independence is used only across the $N$ launches for each indicator.

For completeness, averages of $1,z,z^2$ on three unit intervals centered
at $-1,0,1$ give the invertible matrix with rows $(1,i,i^2+1/12)$.
Its tensor square recovers a coordinatewise quadratic polynomial from
nine cell averages. On scale $h$, Taylor remainder of order $2+\beta$
gives derivative error $Ch^{2+\beta-j}$ at order $j\le2$.
A probability error $\varepsilon$ gives average error $\varepsilon h^{-2}$,
and differentiation costs a further $h^{-j}$. Patch the local polynomials
with a bounded-overlap smooth partition of unity whose derivatives of
order $j$ are $O(h^{-j})$. Leibniz's rule gives the same orders on the
smaller rectangle. At $j=2$, substitution of \eqref{eq:cell-error}
proves \eqref{eq:biased-rate}. No $C^2$ bound on the raw difference
$f_L-f$ is assumed; total variation was first converted to cell errors.

Under \eqref{eq:box-bandwidth}, $\tau_L\le C h^{\beta+4}$.
The bias and square-root terms balance at $r_N^{\beta/(2\beta+6)}$;
the linear Bernstein term is smaller. This proves the last assertion.
\end{proof}

A fixed pilot can select the two active absolute times. Uniform geometry
gives $cd^2\le P(T_0+d)\le Cd^2$ on a common right collar and zero
probability before onset. A fixed fine grid on the known onset bracket,
a fixed positive threshold and a collar after the upper bracket endpoint
place a reported crossing within a small fixed positive distance of $T_0$.
For sufficiently small $\tau_L+\sqrt{r_N}$ the same remains true with
sampling error. Adding two fixed positive offsets then places both
histograms on a common strictly active rectangle with separated times.
The pilot only locates an active interval; it is not an arbitrarily accurate
onset estimator. Fresh preparations for the two selected times justify
conditioning on the pilot. There are a prior-dependent bounded number
of pilot times per record, and their failures can be included in the
constant in Lemma~\ref{lem:biased-histogram}.

\section{A growing record catalogue and its discovery cost}
\label{sec:catalogue}
A positive certificate by itself does not acquire data. This section
specifies the producer access used in our acquisition theorem. The
statistical reconstruction in the next section sees only the records,
not the producer's geometric workspace. We give both an abstract arrival
condition and a concrete realization using the geometric primitives of
the retained scanner. We do not replace those primitives by unmarked
trajectory observations without a proof.

\subsection{A small saturated subcollection}
A \emph{saturated witness} is a finite set of genuine records whose
observed graph is connected, covers every type, and has cycle subgroup
$\Lambda$. It is used in the proof of discovery, not given to the learner.
Recall $B_0,Q$ from \eqref{eq:denominator-bound}.

\begin{lemma}[Witness size without a primitive pair]\label{lem:small-witness}
Every covering saturated record graph on the uniform class contains a
covering saturated subcollection of at most
\begin{equation}\label{eq:witness-size}
                  w_0=r_0+1+\lfloor\log_2 Q\rfloor
\end{equation}
records. The assertion does not require any pair of recorded cycles to
be a basis of $\Lambda$.
\end{lemma}
\begin{proof}
Keep a spanning tree, using at most $r_0-1$ records. Its fundamental
cycle displacements generate $\Lambda$. Choose two independent such
displacements and keep their two non-tree records. Their subgroup has
index $n\le Q$ in $\Lambda$. If the currently kept subgroup is proper,
some remaining fundamental cycle is outside it. Adding this cycle strictly
enlarges the subgroup; its index in $\Lambda$ becomes a proper divisor
of the preceding index, hence is at most half of it. After at most
$\lfloor\log_2 Q\rfloor$ such additions its index is one. A fundamental
cycle needs only its non-tree record and the already retained tree, so
the stated bound counts records, not the lengths of their closed walks.
\end{proof}

\begin{definition}[Predictable genuine-record producer]\label{def:producer}
Before validation at epoch $k$, a producer makes one proposal, using all
previous observations and fresh proposal randomness if desired. It returns
either no record or a replayable genuine record satisfying the local
physical bounds. The cumulative list $\mathcal L_k$ has $J_k\le k$
entries and retains all previous entries, including repetitions. A record
includes its periodic evaluation gate and an onset bracket; it need not
supply the density functions themselves. Histogram preparations at epoch
$k$ are fresh, conditional on the proposal history.

For a fixed unknown saturated witness with $w$ record types, let $H_\ell$
be the set of proposals yielding an instance of witness record $\ell$,
with any admissible local sign and window choices. We say the producer
has discovery hazards $p_{k\ell}$ if, conditionally on the past and while
$\ell$ has not appeared, its next proposal is in $H_\ell$ with probability
at least the deterministic number $p_{k\ell}\in[0,1]$.
The sets $H_\ell$, the witness and these bounds are not inputs to the
stopping test. They are conditions for its termination and cost bounds.
\end{definition}

\begin{lemma}[Discovery under adaptive arrivals]\label{lem:arrival}
For a producer as above, let $K_{\rm cov}$ be the first epoch by which
all witness records have appeared. Then
\begin{equation}\label{eq:arrival-tail}
 \Prob\{K_{\rm cov}>k\}\le
       \sum_{\ell=1}^w\exp\left(-\sum_{j=1}^kp_{j\ell}\right).
\end{equation}
If every cumulative hazard diverges, $K_{\rm cov}<\infty$ almost surely.
If $p_{k\ell}\ge p>0$, the bound is $w e^{-pk}$.
\end{lemma}
\begin{proof}
On the event that record $\ell$ has not appeared before proposal $j$,
its conditional probability of remaining absent is at most $1-p_{j\ell}$.
Iterated conditioning bounds its absence probability by
$\prod_{j\le k}(1-p_{j\ell})\le\exp(-\sum_{j\le k}p_{j\ell})$.
A union bound proves \eqref{eq:arrival-tail}. For divergent hazards the
probability of permanent absence is zero for each of the finitely many
witness records. Independence of different proposals is not required.
\end{proof}

\subsection{A randomized realization of the geometric scanner}
We specify the access precisely. On a bounded laboratory disk a geometric
oracle lists the physical bodies meeting it, represents their complete
boundaries, and tests shortest-pair clearance against all listed blockers.
For a retained clear normal pair it installs the local arclength return
gate at every translated occurrence of that pair and supplies a local
physical chart with verified margins. Translation classes may be recognized
from the congruence of complete pair images on the asymmetric
shape-separated class; no period vectors are required. This is geometric
access, not inference from collision counts. The reconstruction algorithm
receives the installed gates and their margins, not these boundary maps.
The oracle's work is counted as proposal work separately from phase
preparations. No bit complexity for exact analytic-body operations is asserted.

Let $\rho_* =\sup_x\dist(x,\mathcal O)$. A fundamental parallelogram
of diameter at most $b$ and any one obstacle point imply $\rho_*\le b$:
translate the point so that its displacement from $x$ is represented in
that parallelogram. Choose $R>2b$ and a midpoint aperture radius $B>b$.
Put $H=B+R/2$ and list bodies meeting the disk of radius $H$, obtaining
complete bodies in its padding by $D_0$. For distinct bodies separated
by at least $d_0$, a packing bound is
\begin{equation}\label{eq:body-count}
 M\le M_0:=\left\lceil\{1+2(H+D_0)/d_0\}^2\right\rceil.
\end{equation}
For example choose one point in each body in the padded disk; the open
$d_0/2$ disks about these points are disjoint. An epoch chooses uniformly
an ordered pair of distinct listed bodies. It retains the corresponding
shortest bridge if its gap is less than $R$, its midpoint lies in the
aperture, and its verified physical margins are admissible. Otherwise the
proposal returns no record. Previously obtained records remain available.

\begin{proposition}[Finite discovery without a supplied record list]
\label{prop:geometric-producer}
Suppose the clear bridges of gap less than $R$ have the local margins
required above. The randomized producer contains a saturated witness
with $w\le w_0$, and each of its witness records has conditional proposal
probability at least
\begin{equation}\label{eq:producer-p}
                    p=\{M_0(M_0-1)\}^{-1}.
\end{equation}
Thus it satisfies \eqref{eq:arrival-tail} with this $p$. The body list,
the selected witness, the period vectors and the complete record catalogue
are not supplied to the statistical learner. The guarantee uses the
geometric oracle just specified; it is not a no-geometric-access theorem.
\end{proposition}
\begin{proof}
We first show that the lifted graph of all clear shortest bridges of gap
less than $R$ is connected. Choose $\rho_*<a<R/2$. Open $a$-neighborhoods
of the bodies cover the connected plane; their intersection graph is
connected. Each intersecting pair has gap less than $R$. If its shortest
segment meets another body, the gaps to that intermediate body from each
endpoint body are strictly smaller. There are only finitely many pair-gap
values below $R$ modulo periods: only finitely many lattice translations
bring two bounded representatives this close. Induction over these values
replaces every blocked pair by a path of clear smaller-gap pairs. Hence
the clear graph is connected. Its finite quotient covers all obstacle
types. A lifted path from one copy to each of its lattice translates
shows that the closed-walk translations generate every element of $\Lambda$.

Every midpoint translation class has a representative in a fundamental
parallelogram lying within distance $b$ of the aperture center; the actual
parallelogram is not used by the scanner. Thus the aperture contains a
representative of each required bridge. Its endpoint bodies and every
possible blocker meet the disk of radius $H$. The finite listed pair set
therefore contains a covering saturated collection of admitted bridges.
Lemma~\ref{lem:small-witness} reduces it to at most $w_0$ records.
Each has at least one ordered pair representative, selected with probability
$1/(M(M-1))\ge p$. Independent uniform pair choices realize the hazards;
additional adaptive proposals are allowed if this lower probability is
preserved. Apply Lemma~\ref{lem:arrival}.
\end{proof}

A proposal primitive that already returns complete body boundaries is a
strong acquisition resource. Proposition~\ref{prop:geometric-producer}
does not count that information as if it were learned from noisy endpoint
histograms. Its purpose is to close the arrival and stopping analysis of
the previously justified scanner, with an explicit discovery cost and
without requiring that its full list be completed before validation begins.
A different producer may replace it if its genuine-branch validity and
hazards are proved in its own observation model.

\subsection{Countability, effectiveness and exhaustion}
For a fixed periodic table, physical obstacle copies form a countable
locally finite family. Pairs of copies are countable, each disjoint strictly
convex pair has a unique shortest segment, and admissible choices of strictly
active windows may be restricted to rational times in open intervals.
A fixed rational rule shrinking contact rectangles inside a valid chart
then gives a countable sufficient collection of record descriptors. The
underlying density functions still vary in an infinite-dimensional space;
we are not enumerating all analytic functions by finite strings.

With the geometric oracle, increasing integer apertures, bounded-range
pair lists and rational window refinements give an enumeration relative
to that oracle. Without the oracle or a separately proved record producer,
countability alone does not make the collection effectively enumerable.
Exact comparison of arbitrary real-analytic boundaries has no Turing
complexity consequence here. The rational/Hermite stage, once its finite
inputs are separated, is an ordinary finite integer computation.

\begin{proposition}[Pointwise discovery by local-margin exhaustion]
\label{prop:exhaustion}
Keep the global analytic, area, shape and separation priors, but allow
record-specific positive local margins and finite onset brackets certified
by the geometric producer. Suppose a finite covering saturated set of
strictly clear records exists. There is an oracle-relative growing-list
experiment that stops almost surely at every prescribed positive accuracy,
with simultaneous error probability at most $\delta$. No common positive
local margin for the unknown witness need be supplied. This pointwise
assertion gives no uniform stopping rate over the exhaustion levels.
\end{proposition}
\begin{proof}
At level $j$ use a finite aperture and gap cutoff, bounded onset brackets
and time windows, and local margins at least $2^{-j}$, with corresponding
finite chart bounds. Choose these constraints nested so their union
contains every strictly valid finite record. Introduce one level per epoch.
Give each introduced level $j$ proposal probability $2^{-j}$, assigning
the unused probability to a null proposal. Within a level choose uniformly
among its finitely many candidate pairs and admissible descriptors from
the finite refinement used there. Each of the finitely many witness
records, after choosing one admissible rational descriptor, belongs to
some fixed finite level. Thus it has a fixed positive proposal probability
at every later epoch, and is discovered almost surely.

At epoch $k$ validate the records admitted at the first $k$ levels. This
is a finite union with positive minimum margins and finite maximum bounds;
Section~\ref{sec:uniform} applies with known level-dependent constants.
Choose finite sample sizes and launch squares so that the simultaneous
\emph{table} error is at most $2^{-k}$ and all arithmetic and defect
thresholds hold, using failure budget $\delta/[8(k+1)^6]$. This is
possible by Lemma~\ref{lem:biased-histogram}, even if the level-dependent
constants deteriorate. No comparison of powers with changing constants
is used. Keep all admitted records. The first Borel--Cantelli lemma gives
only finitely many inaccurate validation epochs. Once the finite witness
has appeared, a later accurate epoch with $2^{-k}\le\xi$ must accept.
A summable union bound gives the same simultaneous soundness as
Theorem~\ref{thm:sequential}. The resulting sample sizes need not follow
its polynomial schedule; hence its uniform tail and expected-cost bound
are not asserted for this exhaustion.
\end{proof}

\section{Anytime certification of a growing catalogue}
\label{sec:sequential}
The data list is now random, unknown in advance and chosen using earlier
observations. Repeatedly applying a fixed-list confidence statement at
one fixed error probability would not justify stopping. We allocate
summable error probabilities over epochs and, within each epoch, estimate
all retained records using fresh preparations. Once their simultaneous
accuracy is known, the deterministic certificate already covers every
component and every subcollection. It is not necessary to allocate an
additional error probability to each graph or possible stopping choice.

The statistical device is elementary probability spending, not a new
confidence-sequence inequality. General time-uniform inference is developed
in \cite{HRMS}. The work here is to carry a uniform event through analytic
continuation, unknown type incidence, rational decoding, volume calibration,
and discovery of a previously unspecified complete record collection.

\subsection{The experiment}
Fix the uniform global and local class of Section~\ref{sec:uniform}, the
preparation bounds of Proposition~\ref{prop:box}, a genuine-record producer,
$0<\delta<1/4$, and a requested table accuracy $\xi>0$. Record validity
is a property of the producer and physical model, not inferred merely
from an empirical fit. Let
\begin{equation}\label{eq:epoch-schedule}
 \delta_k=\frac{\delta}{8(k+1)^6},\qquad
 N_k=(k+1)^2,\qquad
 r_k=\frac{\log\{C_*(k+1)^9/\delta\}}{(k+1)^2}.
\end{equation}
Choose $C_*$ sufficiently large in terms of the fixed class. At epoch
$k$, first request one new proposal and retain its record, if any, together
with all preceding records. Conditional on this history the list has
$J_k\le k$ entries. When $r_k$ is small enough for the common chart,
use the fixed pilot and two histograms per retained record with $N_k$
fresh accepted free preparations per time, bandwidth (rounded to a centered grid within a fixed factor) and launch square
\begin{equation}\label{eq:epoch-geometry}
 h_k=r_k^{1/(2\beta+6)},\qquad
 L_k=\max\{2b,C_Lr_k^{-(\beta+4)/(2\beta+6)}\}.
\end{equation}
Before this small-error regime collect proposals but make no acceptance
claim. The constants $C_*,C_L$ depend on all the stated priors, not on an
unknown witness or the producer hazards.

Fit each estimated density pair in the compact physical pair class as
in Section~\ref{sec:uniform}. Let
\begin{equation}\label{eq:epoch-radii}
 e_k=C_2r_k^{\beta/(2\beta+6)},\qquad
 t_k=C_3e_k^\theta
\end{equation}
be valid uniform bounds, including the fitting factor. Run the
reference-free graph, rational-coordinate and defect tests only when
$e_k<e_0$, all pilot margins hold and $t_k\le\xi$. A rank failure, a
nonunique rational decoding or a failed defect test returns no certificate.
Stop at the first epoch with an accepted component. Denote this stopping
epoch by $T_\xi$.

Every test uses the same physical area estimate and decoded subgroup
as in Section~\ref{sec:uniform}. In particular it never substitutes
$A+S_I$ for a covolume before certification, and never takes an integer
span of noisy real vectors. Windows and grids may depend on the pilot,
but subsequent preparations are fresh. Reversal is fixed within each
record across its windows and epochs; it need not be coordinated between
records.

\Needspace{7\baselineskip}
\begin{theorem}[Anytime soundness and finite discovery]\label{thm:sequential}
For the experiment above the following hold.

\smallskip\noindent
\textup{(i)} With probability at least $1-\delta$, every certificate at
every epoch is correct: its component covers all obstacle orbits and its
decoded period lattice is the full translation group. Its reconstructed
table has simultaneous support-function $C^2$ and period-basis error at
most $t_k$. These assertions hold for data-dependent component selection,
deletions and stopping. Deletions need not preserve termination.

\smallskip\noindent
\textup{(ii)} If the retained producer has a finite saturated witness
whose cumulative hazards diverge, then $T_\xi<\infty$ almost surely.
No witness or complete list is an input. Let $k_\xi$ be a deterministic
epoch after which all error thresholds and $t_k\le\xi$ hold. For
$k\ge k_\xi$,
\begin{equation}\label{eq:stopping-tail}
 \Prob\{T_\xi>k\}\le\delta_k+
       \sum_{\ell=1}^w\exp\left(-\sum_{j=1}^kp_{j\ell}\right).
\end{equation}
For uniform hazards $p_{j\ell}\ge p>0$, this is at most
$\delta_k+w e^{-pk}$; in particular the expected stopping epoch and
the expected total number of launch attempts are finite.

\smallskip\noindent
\textup{(iii)} Up to any epoch $m$, at most $C(m+1)^4$ accepted free
preparations are used. Simultaneously for all $m$, with probability at
least $1-\delta$, the total launch attempts are at most
\begin{equation}\label{eq:sequential-cost}
 C\{(m+1)^4+(m+1)\log((m+1)/\delta)\}.
\end{equation}
There are $m$ record proposals, in addition to any setup work of the
specified producer. These are explicit upper bounds for this schedule,
not minimax or polynomial-time bounds for analytic-body fitting.
\end{theorem}
\begin{proof}
Conditional on the proposal history at epoch $k$, the list is a fixed
finite genuine list with $J_k\le k$. For the histogram and pilot bounds,
\[
 \frac{\log(CJ_kN_k/\delta_k)}{N_k}\le r_k
\]
after enlarging $C_*$. Lemma~\ref{lem:biased-histogram}, the pilot argument
and \eqref{eq:epoch-geometry} give a conditional failure probability at
most $\delta_k$ for simultaneous density accuracy $e_k$. Iterated
conditioning gives the same unconditional bound. Write $E_k$ for this
epoch accuracy event. Since $\sum_k\delta_k<\delta/8$, the event
$E=\bigcap_kE_k$ has probability at least $1-\delta/8$.

On $E_k$, Theorem~\ref{thm:uniform-main} applies to \emph{every} component
of the epoch's list and to each retained subcollection. Its constants
are independent of repetitions and of $J_k$. It recovers the correct
integer data before using the estimated area. An incomplete rank-two
component has true defect at least $g_0$ and estimated defect error below
$g_0/4$, whereas a complete saturated component has defect zero. Therefore
$|\widehat{\mathcal D}|<g_0/2$ separates the two cases. The reconstruction
error is at most $t_k$. These deterministic implications on $E$ prove
(i), including an adaptively chosen stopping epoch. A fitted physical
model is used for stable reconstruction, not as evidence that an arbitrary
input gate is genuine.

The radii $e_k,t_k$ tend to zero. Choosing a sufficiently large $C_*$
makes the envelope $r_k$ decreasing once testing begins, so some finite
$k_\xi$ satisfies the stated thresholds for all later epochs. By
Lemma~\ref{lem:arrival}, all witness records have appeared almost surely
at a finite epoch. Since the list retains them, it thereafter has a
complete saturated component. Also $\sum_k\Prob(E_k^c)<\infty$;
the first Borel--Cantelli lemma implies that only finitely many $E_k$
fail, without any independence assumption between epochs. Thus after
both finite exceptional times and $k_\xi$, the test must accept. This
proves almost-sure stopping. One may define the underlying proposal and
preparation streams beyond their first acceptance for this argument;
no continuing measurements are required from the implemented experiment.

For the tail bound, at epoch $k\ge k_\xi$, if the witness has appeared
and $E_k$ holds, acceptance has occurred by that epoch. Hence
$\{T_\xi>k\}\subset\{K_{\rm cov}>k\}\cup E_k^c$.
Use \eqref{eq:arrival-tail} and $\Prob(E_k^c)\le\delta_k$ to obtain
\eqref{eq:stopping-tail}. Notice that only the current epoch error enters
this tail, not a sum of all preceding error probabilities. For uniform
hazards this tail is a sum of a term $O(k^{-6})$ and an exponentially
decaying term. It has a finite fourth moment.

There are a uniformly bounded number of pilot and histogram times per
record. Thus accepted free preparations at epoch $k$ number at most
$C J_kN_k\le Ck(k+1)^2$. Summing gives $C(m+1)^4$. Apply
Proposition~\ref{prop:box} to all required free acceptances at epoch $k$,
with failure probability $\delta_k$. The square is the same at that
epoch, and the free-launch probability is at least $p_0$. Even if the
record outcomes are evaluated at different times, free-launch acceptance
does not depend on the gate. An upper bound for all needed acceptances
can be drawn in advance. The epoch launch budget is therefore at most
$C\{k(k+1)^2+\log(1/\delta_k)\}$ outside an event of probability
$\delta_k$. A second summable union bound proves \eqref{eq:sequential-cost}
simultaneously for all $m$. Together with $E$ its probability is still
at least $1-\delta$.

Finally, conditional on the history before an epoch is run, its expected
launch cost is at most $Ck(k+1)^2$. Under uniform hazards,
\[
 \E[\hbox{total attempts}]
 \le C\sum_{k\ge1} k(k+1)^2\Prob\{T_\xi\ge k\}<\infty
\]
by \eqref{eq:stopping-tail}. The preasymptotic sum is finite, the
$k^{-6}$ tail is summable after multiplication by $k^3$, and so is the
exponential term. This proves (ii)--(iii).
\end{proof}

\begin{corollary}[A complete noisy scan with a stopping bound]
\label{cor:scanner-discovery}
Under the explicit geometric-oracle and uniform-margin assumptions of
Proposition~\ref{prop:geometric-producer}, Theorem~\ref{thm:sequential}
applies with $w\le w_0$ and $p$ given by \eqref{eq:producer-p}. For
$0<\alpha<1$, by any epoch
\begin{equation}\label{eq:scan-stopping}
 k\ge\max\left\{k_\xi,\left\lceil p^{-1}\log(w_0/\alpha)\right\rceil\right\},
\end{equation}
the procedure has returned a correct whole-table reconstruction of error
at most $\xi$, except on an event of probability at most $\delta+\alpha$.
It need not first complete or identify the full short-channel catalogue.
\end{corollary}
\begin{proof}
The probability that the unknown witness is still absent is at most
$w_0e^{-pk}\le\alpha$. On the simultaneous accuracy event, the first
complete saturated component passes at the next eligible epoch, and
all previous acceptances are correct. Use Theorem~\ref{thm:sequential}.
\end{proof}

\subsection{What the sequential conclusion changes}
There are three distinct assertions. The fixed-list theorem gives a
conditional whole-table error rate; the present theorem gives validity
over an unbounded sequence of adaptively growing lists and a stopping-time
bound under explicit arrival conditions; the retained hidden-area experiment
gives an $N^{-1/2}$ two-point decision scale. The variables, observations
and losses differ, so their exponents are not combined into a single
``sharp recovery rate.''

No certificate asserts that an unsuccessful producer has exhausted the
physical table. Rank failure, a positive defect, unresolved arithmetic,
insufficient precision and an empty list all mean only that the current
data have not certified completion. Conversely, genuine-record validity
and arrival are not manufactured by the stopping proof: they must follow
from the specified producer. Proposition~\ref{prop:geometric-producer}
provides one realization with geometric access, and the margin-exhaustion
argument gives pointwise noisy stopping beyond a fixed common local margin.
Recovery from unmarked trajectories without those access primitives is a
different observation problem and is not asserted here.

\section{Information, related inverse problems and companion results}
\label{sec:comparison}
\subsection{The nearest inverse-geometry comparisons}
Stefanov--Uhlmann--Vasy~\cite[Theorems 1.1 and 1.3]{SUV} recover a
Riemannian metric, in dimension at least three, from local boundary
distances or global lens data under their strict-convexity and foliation
hypotheses, modulo boundary-fixing diffeomorphisms. Here the ambient metric
is Euclidean; the unknowns are two reflecting arcs, obstacle incidence and
a translation lattice. Proposition~\ref{prop:lens} identifies rather than
conceals the local information: two physical endpoint densities are
equivalent to the generating action and absolute phase normalization.
We do not claim the affine quotient or travel-time differentiation as a
new general lens-rigidity principle.

Gurfinkel--Noakes--Stoyanov~\cite[Theorems 1 and 2]{GNS} relate exterior
travelling times to non-trapping flow conjugacy and obtain obstacle
uniqueness under non-positive curvature and strict convexity. Noakes--Stoyanov
\cite{NS} study exterior travelling times and give constructive results
for particular obstacle configurations. Their enclosing-boundary observations
are not the selected three-impact endpoint laws used here. Nor does the
present geometric producer follow from their inverse results. The local
curvature calculation is retained with its full proof in Section~\ref{sec:local}.

The completion identity uses the same absolute normalization in every
record; unseen bodies change that normalization. Santal\'o-type volume
identities provide a related flow-measure perspective \cite{Santalo}.
The preparation result here is different: it compares a finite laboratory
launch distribution with the unknown quotient law by a periodic boundary
layer, before any inverse has determined a cell. Its total-variation
error must be amplified through cell reconstruction as in
\eqref{eq:biased-rate}. It is not an equidistribution theorem for a
single billiard trajectory.

Marked-length results for dispersing billiards, including
De Simoi--Kaloshin--Leguil~\cite{DSKL}, and the enriched marked-length
Sinai-billiard result of Finamore--Leguil~\cite{FL}, use different global
dynamical data. We assert neither equivalence of those invariants to our
endpoint sensor nor recovery from an ordinary count-only record. Analytic
continuation is conditional on its strip and magnitude bounds
\cite{Trefethen}; the statistical rate does not remove that sensitivity.

\subsection{Discovery, probability control and algorithmic scope}
The passage from finite-list confidence to a stopping rule requires control
for every epoch, including the epoch selected by the data. General confidence
sequence theory treats time-uniform validity \cite{HRMS}. Our proof uses
the simpler device of conditional fresh validation and summable failure
budgets, together with a separate record-arrival bound. No new general
martingale inequality is claimed. The particular content is its compatibility
with geometric fitting, exact rational subgroup recovery, a missing-area
gap and a finite-size preparation bias in one stopping theorem.

The geometric producer extends the retained bounded-aperture scanner.
The obstruction-descent argument is related to proximity-graph constructions
\cite{Toussaint,JT}, but uses strict shortening of body gaps rather than
a triangle inequality. Visibility complexes \cite{PV} concern different
geometric primitives. Their complexity bounds do not transfer to exact
analytic-body access or to discovery from unmarked trajectories. Likewise
Hermite and Smith algorithms \cite{KB,Cohen} make the separated integer
stage finite, not the compact physical fitting stage polynomial-time.
The random proposals cost one pair choice and its oracle operations;
phase samples, rejected launches and the required spatial apertures are
separate resources. At epoch $k$ the latter have half-width $L_k$ in
\eqref{eq:epoch-geometry}. A preparation-count bound does not bound
apparatus travel time or representation complexity.

A failed test remains a noncertificate. Its cause may be missing records,
a proper subgroup, insufficient precision or a producer that has not yet
returned useful branches. It never licenses the conclusion that no complete
physical table exists. Conversely, a sound test cannot force an invalid
or permanently uninformative producer to acquire a genuine witness.
The explicitly geometric producer supplies one sufficient mechanism;
other observation models require their own genuine-record and arrival proof.

\subsection{What is preserved and what is not inferred}
The complete preceding article is supplied unchanged as Supplement C.
Its original fixed-list histogram theorem, hidden-area testing proof and
geometric acquisition comparison remain available. Its retained companions
include the finite-index classification, all-cycle gcd construction,
canonical channels, finite apertures and persistence when unused channels
change. The present primary includes the local inverse, exact defect and
reference-free arithmetic directly, so the sequential proof does not depend
on an unavailable historical assertion.

The three statistical conclusions must not be merged. A fixed-list
whole-table upper bound has exponent $\beta\theta/(2\beta+6)$ in its
stated sample parameter. Our sequential theorem adds a discovery tail and
a schedule-dependent cost. The hidden-area subexperiment has a sharp
$N^{-1/2}$ decision scale at fixed windows. It is not a matching lower
bound for record discovery, analytic continuation or whole-table estimation.
No one of these exponents is advertised as a universal minimax result.

The full moment-compressed inverse, symmetric or repeated-shape registered
theorems, smooth physical jet realizations and count-fiber arguments remain
in the retained complete v18 article. The complete smooth relative-law,
Volterra, Abel and auxiliary development remains Supplement S. These are
included companion results, not external publications or deleted history.
In particular, formal count fibers and smooth Taylor agreement are not
reinterpreted as equality of exact analytic count germs. All new unregistered
whole-table conclusions retain their asymmetric shape-separated physical
class and the explicit producer and preparation contracts.

\subsection*{Companion material and assisted analysis}
Supplement C is the complete preceding reference-free certification article,
with its full retained companion tree. Supplement S contains the complete
smooth relative-law development and its auxiliary document. The source map
identifies each entry and its immutable provenance. These are included
materials, not independently published papers. AI-assisted analysis contributed
to the sequential discovery, preparation-bias and witness-size arguments.
The named author is responsible for checking the proofs and references.
Finite diagnostics and compilation provide reproducibility evidence, not
formal proof certification or independent referee acceptance.
\begin{thebibliography}{99}
\raggedright
\bibitem{Cohen}
H. Cohen, \emph{A Course in Computational Algebraic Number Theory},
Graduate Texts in Mathematics, vol.~138, Springer, Berlin, 1993, Chapter~2.
\bibitem{DSKL}
J. De Simoi, V. Kaloshin, and M. Leguil,
\emph{Marked length spectral determination of analytic chaotic billiards
with axial symmetries}, Invent. Math. \textbf{233} (2023), 829--901.
\href{https://doi.org/10.1007/s00222-023-01191-8}{doi:10.1007/s00222-023-01191-8}.
\bibitem{FL}
D. Finamore and M. Leguil,
\emph{A CAT(0)-approach to the marked length spectral rigidity of Sinai
billiards}, arXiv:2510.18983v1 (2025).
\bibitem{GNS}
T. Gurfinkel, L. Noakes, and L. Stoyanov,
\emph{Travelling times in scattering by obstacles in curved space},
J. Differential Equations \textbf{269} (2020), 9508--9530.
\href{https://doi.org/10.1016/j.jde.2020.05.049}{doi:10.1016/j.jde.2020.05.049}.
\bibitem{HRMS}
S. R. Howard, A. Ramdas, J. McAuliffe, and J. Sekhon,
\emph{Time-uniform, nonparametric, nonasymptotic confidence sequences},
Ann. Statist. \textbf{49} (2021), 1055--1080.
\href{https://doi.org/10.1214/20-AOS1991}{doi:10.1214/20-AOS1991}.
\bibitem{JT}
J. W. Jaromczyk and G. T. Toussaint,
\emph{Relative neighborhood graphs and their relatives},
Proc. IEEE \textbf{80} (1992), 1502--1517.
\bibitem{KB}
R. Kannan and A. Bachem,
\emph{Polynomial algorithms for computing the Smith and Hermite normal
forms of an integer matrix}, SIAM J. Comput. \textbf{8} (1979), 499--507.
\href{https://doi.org/10.1137/0208040}{doi:10.1137/0208040}.
\bibitem{NS}
L. Noakes and L. Stoyanov,
\emph{Travelling times in scattering by obstacles},
J. Math. Anal. Appl. \textbf{430} (2015), 703--717.
\href{https://doi.org/10.1016/j.jmaa.2015.05.013}{doi:10.1016/j.jmaa.2015.05.013}.
\bibitem{PV}
M. Pocchiola and G. Vegter, \emph{The visibility complex},
Proc. 9th Annual ACM Symposium on Computational Geometry (1993), 328--337.
\bibitem{SUV}
P. Stefanov, G. Uhlmann, and A. Vasy,
\emph{Local and global boundary rigidity and the geodesic X-ray transform
in the normal gauge}, Ann. of Math. (2) \textbf{194} (2021), 1--95.
\href{https://doi.org/10.4007/annals.2021.194.1.1}{doi:10.4007/annals.2021.194.1.1}.
\bibitem{Santalo}
L. Stoyanov,
\emph{Santal\'o's formula and stability of trapping sets of positive measure},
J. Differential Equations \textbf{263} (2017), 2991--3008.
\href{https://doi.org/10.1016/j.jde.2017.04.019}{doi:10.1016/j.jde.2017.04.019}.
\bibitem{Toussaint}
G. T. Toussaint, \emph{The relative neighbourhood graph of a finite planar
set}, Pattern Recognition \textbf{12} (1980), 261--268.
\bibitem{Trefethen}
L. N. Trefethen,
\emph{Quantifying the ill-conditioning of analytic continuation},
BIT Numer. Math. \textbf{60} (2020), 901--915.
\href{https://doi.org/10.1007/s10543-020-00802-7}{doi:10.1007/s10543-020-00802-7}.
\end{thebibliography}

\end{document}
